% ID: SEL-D0046-OA-E03
% ESTADO: cotejada. Segunda pasada visual y borrador compilado revisados.
% Fuente: Selectividad/sel_2014_matematicas_aplicadas/Matematicas_CCSS_JUN.pdf
% SHA-256: df7375f8183eafb6f45e9c65083e68d528ffabe07d66717b0ccf0f020c5b064f
% Página: 1; opción A; ejercicio original 3.

Una urna, A, contiene siete bolas numeradas del 1 al 7. Otra urna, B, contiene cinco bolas numeradas del 1 al 5. Lanzamos una moneda equilibrada, de forma que si sale cara, extraeremos una bola de la urna A, y, si sale cruz, la extraemos de la urna B.

Calcule las probabilidades de los siguientes sucesos:
\begin{enumerate}[label=\alph*)]
\item \textbf{(0.5 puntos)} ``La bola haya sido extraída de la urna A y el número sea par''.
\item \textbf{(1 punto)} ``El número de la bola extraída sea par''.
\item \textbf{(1 punto)} ``La bola sea de la urna A, si ha salido un número par''.
\end{enumerate}
